\scanpage{36}
\mspage{\handwritten{4}}
% archive sheet number: 35
\uline{finite presentation, and $R_j$ an equivalence relation on $X_j$ defined by an $S_j$-monomorphism $v_j:R_j\handwritten{\longrightarrow}X_j\handwritten{\times}_{S_j}X_j$ such that the morphism $\operatorname{pr}_{1,j}\mathbin{\unsure{o|\circ}} v_j:R_j\handwritten{\longrightarrow}X_j$ obtained from $v_j$ by composition with the first projection $X_j\handwritten{\times}_{S_j}X_j\handwritten{\longrightarrow}X_j$ is flat and of finite presentation. Then, for the quotient sheaf $X_j\handwritten{\times}_{S_j}S/R_j\handwritten{\times}_{S_j}S$ for the topology $\handwritten{\mathbb{T}}$ to be representable, it is necessary and sufficient that there exist $i\handwritten{\in}I$, $i\handwritten{\geqslant}j$, the quotient sheaf $X_j\handwritten{\times}_{S_j}S_i/R_j{}_{S_j}S_i$ for the topology $\handwritten{\mathbb{T}}$ be representable.}

In view of the statements of (EGA IV 8.8.2, 8.8.3, 8.10.5 and 11.2.6) recalled in (10.0), this lemma is a consequence\deleted{\unsure{of the}} of the following result :

\uline{Lemma 10.8.}- \uline{Let $\handwritten{\mathbb{T}}$ be the fppf or fpqc topology ; let $\deleted{\unsure{j\in I,}}$ $X$ be an $S$-prescheme of finite presentation (resp. locally of finite presentation), $R$ an \handwritten{e}quivalence relation on $X$ defined by a monomorphism $v:R\handwritten{\longrightarrow}X\handwritten{\times}_S X$ such that $\operatorname{pr}_1\mathbin{o} v$ is flat and of finite presentation (resp. flat and locally of finite presentation). Then the following conditions are equivalent :}

\uline{(i) The quotient sheaf $X/R$ for the topology $\handwritten{\mathbb{T}}$ is representable.}

\uline{(ii) there exists an $S$-prescheme of finite presentation (resp. locally of finite presentation ) $Y$ and a faithfully flat morphism $p:X\handwritten{\longrightarrow}Y$ such that the diagram}

% Diagram objects and labels typed; connecting arrows handwritten.
\diagram{\(\text{(D)}\)\qquad\begin{tikzcd}[ampersand replacement=\&,column sep=large,row sep=large]
R \handwritten{\arrow[r,"\operatorname{pr}_1\mathbin{o} v"]} \handwritten{\arrow[d,"\operatorname{pr}_2\mathbin{o} v"']} \& X \handwritten{\arrow[d,"p"]} \\
X \handwritten{\arrow[r,"p"']} \& X
\end{tikzcd}}

(where $\operatorname{pr}_1$ and $\operatorname{pr}_2$ are the projections $X\handwritten{\times}_S X\handwritten{\unsure{\substack{\longrightarrow\\[-.6ex]\rightharpoonup}|\rightrightarrows}}X$) \uline{be cartesian.}

Note first that according to (Exp. IV 3.3.2 and 4.4.3), for the sheaf $X/R$ \deleted{\illegible{1}} for the topology $\handwritten{\mathbb{T}}$ to be representable by $Y$, it is necessary and sufficient that the diagram (D) be cartesian and that $p$ be covering for the topology $\handwritten{\mathbb{T}}$.

Let us show that i) implies ii). The hypothesis (i) implies that the diagram (D) is cartesian, hence that $\operatorname{pr}_1\mathbin{o} v$ is obtained from $p$ by base change along $p$, and that $p$ is
